\section{Open-closed TQFTs and the Cardy Condition}
\label{sec:cardy}

We conclude with a discussion of open-closed topological field theories in the B-model and we prove that a condition holds for
Hochschild homology which is equivalent to the Cardy Condition in the
Calabi-Yau case.  Appropriate references for open-closed 2d
topological field theories include Moore-Segal~\cite{MooSeg},
Costello~\cite{Cos} and Lauda-Pfeiffer~\cite{LaudaPfeifer}.

 \subsection{Open-closed 2d TQFTs.}
Consider the open and closed 2-cobordism category $\TwoCob$ whose
objects are oriented, compact one-manifolds --- in other words,
disjoint unions of circles and intervals --- and whose morphisms are
(diffeomorphism classes of) cobordisms-with-corners between the source
and target one-manifolds.  A morphism can be drawn as a vertical
cobordism, from the source at the bottom to the target at the top.  As
well as parts of the boundary being at the top and the bottom, there
will be parts of the boundary in between, corresponding to the fact
that this is a cobordism with corners.  An example is shown below.
\[\pic{cobordism}\]
Disjoint union makes $\TwoCob$ into a symmetric monoidal category and
an open-closed two-dimensional topological quantum field theory (2d
TQFT) is defined to be a symmetric monoidal functor from $\TwoCob$ to
some appropriate symmetric monoidal target category, which we will
take to be the category of vector spaces or the category of graded
vector spaces.  The category $\TwoCob$ has a simple description in
terms of generators and relations which means that there is a
reasonably straight forward classification of open-closed 2d TQFTs up
to equivalence.  This is what we will now describe.  The following
morphisms generate $\TwoCob$ as a symmetric monoidal category.
\[\pic{fig1}\]
We will come back to the relations below.


To specify an open-closed 2d TQFT up to equivalence on objects it
suffices to specify the image $\closed$ of the circle and the image
$\open$ of the interval.  The former is called the space of closed
string states and the latter is called the space of open-string
states.  Using the four \emph{planar} generating morphisms pictured
above, together with the relations between them, it transpires that
$\open$, the space of open-string states is precisely a symmetric, but
not-necessarily commutative Frobenius algebra. This means that it is a
unital algebra with a non-degenerate, symmetric, invariant
inner-product.  It is useful to note here that the inner product is
symmetric because the two surfaces pictured below are diffeomorphic,
however these surfaces are \emph{not} ambient isotopic --- so one
cannot be deformed to the other in three-space while the bottom
boundary is fixed.
\[\pic{syminnerprod1}\stackrel{\text{diffeo}}{=}\pic{syminnerprod2}\]
On the other hand, the first four generating morphisms, along with
their relations, mean that $\closed$, the space of closed string
states, is a \emph{commutative} Frobenius algebra. The last two
morphisms mean that there are maps $i_*\colon \closed\to\open$ and
$i^*\colon \open \to \closed$, and by the relations these are adjoint
with respect to the pairings on these spaces.  Moreover, $i_*$ is an
algebra map, such that its image lies in the centre of $\open$. The
final relation that these must satisfy is the \emph{Cardy Condition}.
In terms of the generators pictured above this is the following
relation:
\[\pic{cardy1}=\pic{cardy2}.\]
Note again that these surfaces are diffeomorphic but not isotopic in
three-space.  In terms of maps, writing $\mu\colon \open\otimes
\open\to \open$ and $\delta\colon\open\to\open\otimes\open$ for the
product and coproduct of the open string state space and writing
$\tau\colon\open\otimes\open\to\open\otimes\open$ for the symmetry in
the target category, the Cardy Condition is the equality of maps from
$\open$ to $\open$:
\[\mu\circ \tau\circ \delta=i_*\circ i^*.\]
We will have reason to use an equivalent condition below.


To summarize, having an open-closed 2d TQFT is equivalent to having
the data of a commutative Frobenius algebra $\closed$, a symmetric
Frobenius algebra $\open$, and an algebra map $i_*\colon \closed \to
\open$ with central image, such that the Cardy Condition is satisfied.

 \subsection{Open-closed 2d TQFTs with D-branes.}
A more interesting model of string theory is obtained when we specify
a set of `boundary conditions' or `D-branes' for the open strings. For
a mathematician this just means a set of labels for the boundary
points of objects.  So fix a set $\Lambda$ of labels, and consider the
category $\TwoCobL$ of open-closed cobordisms such that the objects
are compact, oriented one-manifolds with the boundary points labelled
with elements of the set $\Lambda$, and morphisms having their
internal boundaries labelled compatibly with their boundaries. Here is
an example of a morphism from the union of the circle and the interval
labelled $(B,A)$, to the interval labelled $(B,A)$.
\[\pic{labelledcobordism}\]


Now a $\Lambda$-labelled open-closed TQFT is a symmetric monoidal
functor to some appropriate target category which we will again take
to be the category of vector spaces or the category of graded vector
spaces. Moreover, the category $\TwoCobL$ is similarly generated by
morphisms as listed above, but now they must all be labelled, and the
relations are just labelled versions of the previous relations.  Thus
we can similarly classify $\Lambda$-labelled open-closed TQFTs.  Once
again the image of the circle is a commutative Frobenius algebra,
$\closed$.  However, rather than getting a single vector space $\open$
associated to an interval, we get a vector space $\open_{BA}$
associated to each ordered pair $(B,A)$ of elements of $\Lambda$; so
we do not get a single Frobenius algebra, but rather something which
could be called a `Frobenius algebra with many objects' or a
`Frobenius algebroid', but, for the reason explained below, such a
thing is commonly known as a Calabi-Yau category.  It is a category in
the following sense.  We take the category whose objects are
parametrized by $\Lambda$ and, for $A,B\in \Lambda$, the morphism set
$\Hom(A,B)$ is taken to be $\open_{BA}$ (this is consistent with us
reading diagrams from right to left).  The composition
$\mu_{CBA}\colon\open_{CB}\otimes\open_{BA}\to\open_{CA}$ is given by
the image of the appropriately labelled version of the morphism
pictured.
\[\pic{labelledproduct}\]
The Frobenius or Calabi-Yau part of the structure is a --- possibly
graded --- perfect pairing $\open_{AB}\otimes\open_{BA}\to \k$: the
grading degree of this map is called the dimension of the Calabi-Yau
category.


So to specify a labelled open-closed 2d TQFT it suffices to specify a
commutative Frobenius algebra $\closed$, a Calabi-Yau category $\open$
and an algebra map $i_A\colon \closed\to \open_{AA}$ with central
image, for each object $A$, such that the labelled version of the
Cardy Condition holds.

 \subsection{The open-closed 2d TQFT from a Calabi-Yau
  manifold.}
Associated to a Calabi-Yau manifold $X$ there are two standard 2d TQFTs
coming from string theory, imaginatively named the 
A-model and the B-model:  it is the B-model we will be
interested in here.  In the B-model the boundary conditions are
supposed to be ``generated'' by complex submanifolds of $X$ so the boundary conditions are
taken to be complexes of coherent sheaves on $X$;  the open string
category is then supposed to be the derived category of coherent sheaves on
$X$.  This is indeed a Calabi-Yau category, which is why such
categories are so named:  for each $\cE$ and $\cF$, the requisite pairing
$\Hom^\blob_{\D(X)}(\cE,\cF)\otimes \Hom^\blob_{\D(X)}(\cF,\cE)\to
\k[-\dim X]$ comes from the Serre pairing as a Calabi-Yau
manifold is precisely a manifold with a trivial canonical bundle.


According to the physics, the closed string state space $\closed$ should be $\Hom^\blob_{\D(X\times X)}(\cO_\Delta,\cO_\Delta)$, in other words, the
Hochschild cohomology algebra
 $\HH^\blob(X)$.
%  This agrees with the fact that $\closed$ should be closely related to
%the center of $\open$.  Indeed, Hochschild cohomology is a derived
%version of the notion of center of a ring or a category.
As $X$ is Calabi-Yau, a trivialization of the canonical bundle induces
an isomorphism between Hochschild cohomology and Hochschild homology,
up to a shift. This means that the closed string space $\closed$
has both the cohomological product and the Mukai pairing, and these make  $\closed$ into a
Frobenius algebra.

We need to specify the algebra maps $i_\cE\colon \closed \to
\open_{\cE\cE}$.  These are maps
\[ i_\cE\colon \Hom^\blob(\cO_\Delta,\cO_\Delta)\to
\Hom^\blob(\cE,\cE) \]
which can be given by interpreting $\cE$  as a kernel $\pt\to X$ and taking $i_\cE$ to be convolution with the identity on $\cE$. This is given
diagrammatically on an element
$\phi\in\Hom^\blob(\cO_\Delta,\cO_\Delta)$ as follows.
\[\pic{phiinHH}\quad\mapsto\quad \pic{iEofphiinHH}.\]

At this point it should be noted that $\closed$ is to be thought of as the centre of the category $\open$.  The notion of centre is generalized from algebras to categories by taking the centre of a category to be the natural transformations of the identity functor; however, in a 2-category an appropriate notion of the centre of an object is the set of 2-endomorphisms of the identity morphism on that object.  This means that $\closed$ is the centre, in this sense, of the category $\open$ in the 2-category $\gVar$.

The map going the other way, $i^\cE\colon \Hom^\blob(\cE,\cE)\to
\Hom^\blob(\cO_\Delta,\cO_\Delta)$ is given by taking the trace, namely for
$e\in \Hom_{\D(X)}(\cE,\cE)$ the map is given by
\[\pic{ereallyinHomEE}\quad\mapsto \quad\pic{iEofeinHomEE}.\]
This definition relies on the fact that $X$ is Calabi-Yau, so that the
Serre kernel is, up to a shift, just the identity 1-morphism $\Id_X$.


%
%
%There is an obvious map from closed states to open states labelled by
%$\cE$ given by convolution with the identity:
%\[ i_\cE\colon \Hom^\blob(\cO_\Delta,\cO_\Delta)\to
%\Hom^\blob(\cE,\cE). \]
%Interpreting $\cE$ as an object of $D(\pt\times X)$, this is given
%diagrammatically on an element
%$\phi\in\Hom^\blob(\cO_\Delta,\cO_\Delta)$ as follows.
%\[\pic{phiinHH}\quad\mapsto\quad \pic{iEofphiinHH}.\]
%The dual operation $i^\cE\colon \Hom(\cE,\cE)\to
%\Hom(\cO_\Delta,\cO_\Delta)$ is given by taking the trace, namely for
%$e\in \Hom_{\D(X)}(\cE,\cE)$ the map is given by
%\[\pic{ereallyinHomEE}\quad\mapsto \quad\pic{iEofeinHomEE}.\]
%This trace is again not defined in general; its definition requires
%$X$ to be Calabi-Yau.  In fact, the natural target space for the trace
%is precisely the Hochschild {\em homology} $\HH_\blob(X)$, in a
%fashion similar to the definition of the Chern character.  The target
%of the trace map can then be taken to be $\HH^\blob(X)$, under the
%identification $\HH^\blob(X) \iso \HH_{\dim X - \blob}(X)$.

An argument similar to Proposition~\ref{prop:iotasadjoint} shows that
$i_\cE$ and $i^\cE$ are adjoint.  In order to argue that we indeed
have an open-closed TQFT it remains to show that the Cardy Condition
holds.  In fact, we will prove a more general statement, the Baggy
Cardy Condition.

 \subsection{The Baggy Cardy Condition.}
In the case of a manifold $X$ that is not necessarily Calabi-Yau we
don't have the same coincidence of structure as above: we no longer
have a Frobenius algebra $\HH^\blob(X)$; rather we have an algebra
$\HH^\blob(X)$ and an inner product space $\HH_\blob(X)$.  This means
that we can not formulate the Cardy Condition as it stands.  We now state a
condition which makes sense for an arbitrary, non-Calabi-Yau manifold and
which is equivalent to the Cardy Condition in the  Calabi-Yau case.

\begin{Theorem}
  Suppose that $\open$ is a Calabi-Yau category and $\closed$ is an
  inner product space, such that for each $A\in \open$ there are
  adjoint maps $i^A\colon \open_{AA}\to \closed$ and $i_A\colon
  \closed \to \open_{AA}$.  Then the Cardy Condition
  \[\mu_{BAB}\circ\tau\circ\delta_{ABA}=i_B\circ i^A\]
  is equivalent to the following equality holding for all $a\in
  \open_{AA}$ and $b\in \open_{BB}$, where the map ${}_{a}m_{b}\colon
  \open_{AB}\to \open_{AB}$ is the map obtained by pre-composing with
  $a$ and post-composing with $b$:
  \[\left\langle i^B{-},i^A{-}\right\rangle_\closed=
  \Tr{}_{-}m_{-}.\]
\end{Theorem}

\begin{Proof}
  The first thing to do is examine the left-hand side of the Cardy
  Condition.  As $\open$ is a Calabi-Yau category there is the
  following equality of morphisms $\open_{AA}\to \open_{BB}$.
  \[\pic{cardyA}=\pic{cardyB}\]
  Note that this does not require any reference to $\closed$, but it
  does fundamentally require the symmetry of the inner product.  This
  is reflected in the fact that the surfaces underlying the above
  pictures are diffeomorphic but not ambient isotopic.

  This means that the Cardy Condition is equivalent to the following
  equality.
  \[\pic{cardyC}=\pic{cardyD}\]
  By the non-degeneracy of the inner product on $\open_{BB}$ this is
  equivalent to the equality of two maps $\open_{BB}\otimes
  \open_{AA}\to\k$ which are drawn as follows.
  \[\pic{cardyE}=\pic{cardyF}\]
  The right-hand side is instantly identifiable as $\left\langle
    i^B{-},i^A{-}\right\rangle_\closed$.  The left-hand side is
  identifiable as the trace of the triple composition map
  $\open_{BB}\otimes\open_{BA}\otimes\open_{AA}\to \open_{BA}$ which
  gives the required result.
\end{Proof}

We can now show that the alternative condition given in the above
theorem holds for the derived category and Hochschild homology of
\emph{any} space: in particular, the Cardy Condition holds for
Calabi-Yau spaces.

\begin{Theorem}[The Baggy Cardy Condition]
\label{thm:cardy}
Let $X$ be a space, let $\cE$ and $\cF$ be objects in $\D(X)$ and
consider morphisms
\[ e\in \Hom_{\D(X)}(\cE, \cE) \quad\text{and}\quad f\in
\Hom_{\D(X)}(\cF, \cF). \]
Define the operator
\[ {}_f m_e \colon\Hom^\blob_{\D(X)}(\cE, \cF) \ra
\Hom^\blob_{\D(X)}(\cE, \cF) \]
to be post-composition by $f$ and pre-composition by $e$.  Then we
have
\[ \Tr {}_f m_e=\MP{\iota^\cE(e)}{\iota^\cF(f)}, \]
where $\iota^\cE$, $\iota^\cF$ are the maps defined in
Section~\ref{subsec:iotae}, and $\Tr$ denotes the (super) trace.
\end{Theorem}

\begin{Proof}
  The proof is very similar to the proof of the
  Semi-Hirzebruch-Rie\-mann-Roch Theorem (Theorem~\ref{thm:hrr}).  The first thing
  to observe is that $\Hom(\cE,\cF)\cong \cE^\chk\circ \cF$ and that
  ${}_f m_e$ is just $e^\chk\circ f$. However, we have
  $\tau_R(e)=\SerK_\pt\circ e^\chk$ and $\SerK_\pt$ is trivial so
  $e^\chk=\tau_R(e)$.  Putting this together with the invariance of
  the Serre trace under the partial trace we get the following
  sequence, and hence the required result.
  \begin{align*}
    \Tr {}_f m_e&=\Tr\left(\pic{baggy1}\right)
    =\Tr\left( \pic{baggy2}\right)\\
    &=\Tr\left( \pic{baggy3}\right)
    =\Tr\left( \pic{baggy4}\right)\\
    &=\MPbig{\pic{baggy5a}}{\pic{baggy5b}}
    =\MP{\iota^\cE\!(e)}{\iota^\cF\!(f)}.\qedhere
  \end{align*}
\end{Proof}
\medskip

\noindent
Observe that the Semi-Hirzebruch-Riemann-Roch Theorem is a direct
consequence of the Baggy Cardy Condition, with $e=\id_\cE$,
$f=\id_\cF$.
